% Part: normal-modal-logic
% Chapter: syntax-and-semantics
% Section: normal-modal-logics

\documentclass[../../../include/open-logic-section]{subfiles}

\begin{document}

\olfileid{nml}{syn}{nor}

\olsection{સામાન્ય મોડલ તર્કો}

મોડલ તર્કને યોગ્ય રીતે સુવ્યવસ્થિત
!!{formula}sનો ગણ માની શકાય.
કેટલાક મોડલ તર્કો ખાસ રસના
છે: તેમને રસપ્રદ અર્થઘટનો અને
ઉપયોગો છે, પરંતુ મુખ્યત્વે તેઓ
સૈદ્ધાંતિક રીતે સારી રીતે
સમજાયેલા છે. તેમને કહેવાતા
સામાન્ય મોડલ તર્કો કહે છે.

\begin{defn}\ollabel{defn:modal-logic}
\emph{મોડલ તર્ક}~$\Log L$
એ મૂળભૂત મોડલ ભાષાનાં
!!{formula}sનો એવો ગણ છે જે
\begin{enumerate}
\item બધી વિધાનાત્મક પુનરુક્તિઓ
  ધરાવે છે,
\item એકરૂપ પ્રતિસ્થાપન હેઠળ
  બંધ છે, અને
\item મોડસ પોનેન્સ હેઠળ બંધ છે:
  જો $!A$ અને $(!A \lif
  !B) \in \Log L$ હોય,
  તો $!B \in \Log L$ પણ.
\end{enumerate}
\end{defn}

\begin{ex}
ખૂબ રસપ્રદ ન હોવા છતાં
આ વ્યાખ્યા સંતોષતા,
અને તેથી મોડલ તર્કો
ગણાતા, !!{formula}sના
કેટલાક ગણ આ છે:
\begin{enumerate}
\item બધી પુનરુક્તિઓનો ગણ;
\item બધા !!{formula}sનો ગણ
  (અસુસંગત મોડલ તર્ક).
\end{enumerate}
\end{ex}

\begin{defn}\label{defn:normal-modal-logic}
મોડલ તર્ક~$\Log L$
\emph{સામાન્ય} ત્યારે અને
માત્ર ત્યારે જ જ્યારે
\begin{enumerate}
\item તે નીચેનાં બંને
  સૂત્રો ધરાવે:
\begin{align}
\tag{K} \Box(p \lif q) & \lif (\Box p \lif \Box q) \\
\tag{Dual} \Diamond p & \liff \lnot \Box \lnot p
\end{align}
\item અને તે આવશ્યકીકરણ
  હેઠળ બંધ હોય: જો
  $!A \in
  \Log L$ હોય,
  તો $\Box !A \in \Log L$
  પણ હોય.
\end{enumerate}
\end{defn}

મોડલ તર્કો વ્યાખ્યાયિત કરવાની
બે મુખ્ય રીતો આપણે જોઈશું.
એક રીત સાબિતી-સૈદ્ધાંતિક
છે: અમુક !!{derivation}
પ્રણાલીઓમાં નિષ્પન્ન થઈ
શકતાં !!{formula}sનો ગણ
લેવો. બીજી રીત
નિદર્શ-સૈદ્ધાંતિક છે:
નિદર્શોના અમુક વર્ગોમાં
પ્રામાણ્ય ધરાવતાં
!!{formula}sના ગણ લેવા.
આ સામાન્ય વિધાનાત્મક
તર્કની, અને તે બાબતે
પ્રથમ-ક્રમ તર્કની પણ,
પરિસ્થિતિને મળતું આવે
છે. ત્યાં પણ
!!{formula}sનો ગણ
બે રીતે અલગ તારવી
શકીએ: બધી સત્યમૂલ્ય
નિયુક્તિઓમાં સાચી
પુનરુક્તિઓના ગણ તરીકે,
અથવા કોઈ !!{derivation}
પ્રણાલીમાં બધાં
!!{derivable}
!!{formula}sના ગણ
તરીકે. સામાન્ય
મોડલ તર્કો માટે
આ બે રીતો
સંબંધાત્મક નિદર્શો
અને સ્વયંસિદ્ધિમૂલક
(હિલ્બર્ટ-પ્રકારની)
સાબિતી પ્રણાલીઓ
વડે શક્ય બને છે.
\end{document}
